\documentclass[12pt]{article}
\usepackage[utf8]{inputenc}
\usepackage[T1]{fontenc}
\usepackage{hyperref}
\title{Kioko Language Resources for K-12 Curriculum Development}
\author{Compiled by Claude Sonnet 5 (claude-sonnet-5) for Andr\'as Kornai}
\date{}
\begin{document}
\maketitle

\section{Introduction}
Kioko (ISO 639-3 ues; Glottocode kiok1239) is an Austronesian language of Buton Island, Indonesia, in the Celebic and Muna--Buton group. Speaker estimates vary widely: about 1,000 (Verbix), 5,970 (ASJP) and about 8,500 (Joshua Project). Some sources treat Kioko as a dialect of Pancana. The searches found no educational resources at all, so this report is very short.

\section{Primary Educational Resources}
I found no primers, textbooks or school programmes for Kioko.

\section{Linguistic and Reference Materials}
\begin{itemize}
\item The Austronesian Basic Vocabulary Database has a Kioko word list and notes the relationship to Pancana. \url{https://abvd.eva.mpg.de/austronesian/language.php?id=969}
\item The Max Planck Institute's numerals page for Kioko says more word lists are needed. \url{https://lingweb.eva.mpg.de/channumerals/Kioko.htm}
\item The Glottolog record for Kioko. \url{https://glottolog.org/resource/languoid/id/kiok1239.bigmap.html}
\item The Wikipedia article on the neighbouring Pancana language. \url{https://en.wikipedia.org/wiki/Pancana_language}
\end{itemize}

\section{Religious and Community Resources}
\begin{itemize}
\item Joshua Project has a profile of the Kioko-Kambowa people. \url{https://joshuaproject.net/people_groups/12704}
\end{itemize}

\section{Digital and Audio Resources}
I found none.

\section{Curriculum Development Context}
There is almost no documentation to build on. Any curriculum would first need fieldwork, a decision on whether Kioko is a distinct language or a Pancana dialect, and an orthography, and would probably have to use Indonesian as the language of instruction for now.

\end{document}